\documentclass[a4paper,12pt]{article}
\usepackage[UTF8]{ctex}
\usepackage{amsmath, amssymb, amsthm}
\usepackage{geometry}
\usepackage{graphicx}
\usepackage{enumitem}
\usepackage{fancyhdr}
\usepackage{hyperref}

\geometry{left=2.5cm, right=2.5cm, top=2.5cm, bottom=2.5cm}

\title{动力系统期末试卷解答}
\author{GitHub Copilot}
\date{\today}

\newtheorem{theorem}{定理}
\newtheorem{lemma}{引理}
\newtheorem{definition}{定义}
\newtheorem{problem}{题目}
\newenvironment{solution}{\begin{proof}[解答]}{\end{proof}}

\begin{document}

\maketitle

\section*{一、判断题 (20分)}

\begin{enumerate}
    \item \textbf{题目}：圆周 $S^1$ 上的连续映射其拓扑熵等于零。
    
    \textbf{解答}：\textbf{错误}。
    
    \textbf{理由}：考虑圆周 $S^1 = \mathbb{R}/\mathbb{Z}$ 上的倍增映射 $T(x) = 2x \pmod 1$（即 $z \mapsto z^2$ 在复平面单位圆上的表示）。该映射的拓扑熵为 $h(T) = \ln 2 > 0$。只有圆周上的同胚映射（如旋转）或度数为 0, 1, -1 的某些映射拓扑熵可能为 0，但一般连续映射（特别是度数 $|d| > 1$ 的映射）拓扑熵不为零。一般地，对于 $S^1$ 上的度数为 $d$ 的映射，其拓扑熵 $h(T) \ge \ln |d|$。

    \item \textbf{题目}：对于几乎所有的实数 $x$，其连分数中的数字序列 $\{a_k(x): k \in \mathbb{N}\}$ 的平均值 $(a_1(x) + \dots + a_n(x))/n$ 收敛。
    
    \textbf{解答}：\textbf{正确}。
    
    \textbf{理由}：这是遍历定理的一个应用。连分数展开对应于高斯映射 $G(x) = \{1/x\}$ 在 $(0,1)$ 上的动力系统。高斯映射关于高斯测度 $\mu(A) = \frac{1}{\ln 2} \int_A \frac{1}{1+x} dx$ 是遍历的。数字 $a_k(x)$ 是关于该动力系统的可积函数（实际上 $a_1(x) = \lfloor 1/x \rfloor$ 的积分是发散的，题目中的平均值是否收敛取决于 $a_k$ 的期望是否存在）。
    
    \textit{更正与深入分析}：注意 $a_1(x) = k$ 当且仅当 $x \in (1/(k+1), 1/k]$。
    $\int a_1(x) d\mu(x) \approx \sum k \cdot \frac{1}{k^2} \approx \sum \frac{1}{k} = \infty$。
    因为期望无穷大，根据 Birkhoff 遍历定理，如果积分发散，平均值通常趋于无穷大。
    但是，题目表述为“收敛”，通常指收敛到一个有限数。如果期望无穷大，则该平均值几乎处处发散到 $+\infty$。
    
    \textit{重新审视题目}：题目写的是“平均值...收敛”。
    Khinchin 常数涉及的是几何平均值 $(a_1 \dots a_n)^{1/n}$，它几乎处处收敛。
    算术平均值 $\frac{1}{n} \sum a_k$ 实际上对于几乎所有 $x$ 是发散到 $+\infty$ 的。
    因此，如果题目意思是“收敛到一个有限数”，则该命题为\textbf{错误}。
    如果题目意思是“有极限”（允许 $+\infty$），则是正确的。
    但在数学分析语境下，“收敛”通常指收敛于有限数。
    
    \textbf{结论}：基于标准测度论结果，连分数部分 $a_k$ 的期望无穷大，算术平均发散。故判断为\textbf{错误}。（注：若题目指几何平均则为真，但符号为算术平均）。

    \item \textbf{题目}：设 $X = \mathbb{R}^n/\mathbb{Z}^n$ 为 $n$ 维环面，$m_{Leb}$ 是 $X$ 上的 Lebesgue 测度，$A$ 是 $n$ 阶整数满秩方阵，则 $m_{Leb}$ 是连续映射 $\sigma_A: \mathbb{R}^n/\mathbb{Z}^n \to \mathbb{R}^n/\mathbb{Z}^n, x+\mathbb{Z}^n \mapsto A\cdot x + \mathbb{Z}^n$ 的最大熵测度。
    
    \textbf{解答}：\textbf{正确}。
    
    \textbf{理由}：对于环面自同态 $\sigma_A$，其拓扑熵由 $h(\sigma_A) = \sum_{|\lambda_i|>1} \ln |\lambda_i|$ 给出，其中 $\lambda_i$ 是 $A$ 的特征值。Lebesgue 测度（即 Haar 测度）是 $\sigma_A$ 的不变测度，且其测度熵 $h_{m_{Leb}}(\sigma_A)$ 等于拓扑熵（这是著名的 Pesin 公式或相关结果在环面同态上的特例）。由于变分原理 $h_{top}(T) = \sup_\mu h_\mu(T)$，且 Haar 测度达到了这个上界，因此它是最大熵测度。

    \item \textbf{题目}：设 $X$ 是紧致度量空间，$T: X \to X$ 连续映射。如果 $|M(X, T)| = 1$，那么 $M(X, T)$ 中的唯一测度是遍历测度。
    
    \textbf{解答}：\textbf{正确}。
    
    \textbf{理由}：$M(X, T)$ 表示 $T$-不变概率测度的集合，它是一个非空凸紧集。遍历测度集 $\mathcal{E}(T)$ 恰好是 $M(X, T)$ 的极点集。如果 $M(X, T)$ 只有一个元素 $\mu$，即 $|M(X, T)| = 1$（唯一遍历性），那么这个唯一的元素 $\mu$ 必须是极点（因为凸集由其极点生成），因此 $\mu$ 是遍历的。
\end{enumerate}

\section*{二、计算与证明题}

\begin{problem}[10分]
    设 $X = \mathbb{R}^3/\mathbb{Z}^3$ 三维环面，$m_{Leb}$ 是 $X$ 上的 Lebesgue 测度。
    $$
    A = \begin{pmatrix} 1 & 1 & 1 \\ 1 & 1 & 2 \\ 1 & 2 & 5 \end{pmatrix}
    $$
    试判断映射 $\sigma_A: \mathbb{R}^3/\mathbb{Z}^3 \to \mathbb{R}^3/\mathbb{Z}^3, \sigma_A(x+\mathbb{Z}^3) = Ax + \mathbb{Z}^3$ 关于 $m_{Leb}$ 是否为保测遍历映射。
\end{problem}

\begin{solution}
    \textbf{1. 保测性判断：}
    线性环面自同态 $\sigma_A$ 保 Lebesgue 测度的充要条件是 $|\det(A)| = 1$（如果是整数矩阵，这保证了它是自同构；如果是自同态，只要满射且局部保测，通常也保 Haar 测度，但题目通常隐含 $A \in GL(n, \mathbb{Z})$ 或 $SL(n, \mathbb{Z})$ 语境）。
    计算 $A$ 的行列式：
    $$
    \det(A) = 1 \begin{vmatrix} 1 & 2 \\ 2 & 5 \end{vmatrix} - 1 \begin{vmatrix} 1 & 2 \\ 1 & 5 \end{vmatrix} + 1 \begin{vmatrix} 1 & 1 \\ 1 & 2 \end{vmatrix}
    $$
    $$
    = 1(5-4) - 1(5-2) + 1(2-1) = 1 - 3 + 1 = -1
    $$
    由于 $|\det(A)| = 1$，且 $A$ 是整数矩阵，$\sigma_A$ 是环面自同构，因此它是\textbf{保测}的（保持 Lebesgue 测度不变）。

    \textbf{2. 遍历性判断：}
    环面自同构 $\sigma_A$ 是遍历的，当且仅当矩阵 $A$ 的没有任何特征值是单位根（即 $\lambda^k \neq 1$ 对所有 $k \ge 1$）。
    计算 $A$ 的特征多项式 $P(\lambda) = \det(\lambda I - A)$。
    $$
    P(\lambda) = \lambda^3 - \text{tr}(A)\lambda^2 + (\text{二阶主子式之和})\lambda - \det(A)
    $$
    $\text{tr}(A) = 1 + 1 + 5 = 7$.
    二阶主子式之和：
    $M_{11} = \begin{vmatrix} 1 & 2 \\ 2 & 5 \end{vmatrix} = 1$,
    $M_{22} = \begin{vmatrix} 1 & 1 \\ 1 & 5 \end{vmatrix} = 4$,
    $M_{33} = \begin{vmatrix} 1 & 1 \\ 1 & 1 \end{vmatrix} = 0$.
    Sum $= 5$.
    $\det(A) = -1$.
    $$
    P(\lambda) = \lambda^3 - 7\lambda^2 + 5\lambda + 1
    $$
    我们需要检查 $P(\lambda)$ 是否有根为单位根。
    代入 $\lambda = 1$: $P(1) = 1 - 7 + 5 + 1 = 0$.
    因此，$\lambda = 1$ 是 $A$ 的一个特征值。
    由于 $1$ 是单位根，根据遍历性判别准则，$\sigma_A$ \textbf{不是遍历的}。

    \textbf{结论}：$\sigma_A$ 是保测映射，但不是遍历映射。
\end{solution}

\begin{problem}[10分]
    设 $X = \mathbb{R}^2/\mathbb{Z}^2$ 二维环面，$m_{Leb}$ 是 $X$ 上的 Lebesgue 测度，$\alpha \in \mathbb{R} \setminus \mathbb{Q}$。映射
    $$
    T: X \to X, \quad T\begin{pmatrix} x \\ y \end{pmatrix} = \begin{pmatrix} 5 & 2 \\ 2 & 1 \end{pmatrix} \begin{pmatrix} x \\ y \end{pmatrix} + \begin{pmatrix} \alpha \\ 0 \end{pmatrix} \pmod 1
    $$
    试判断映射 $T$ 关于 $m_{Leb}$ 是否为保测遍历映射。
\end{problem}

\begin{solution}
    记 $A = \begin{pmatrix} 5 & 2 \\ 2 & 1 \end{pmatrix}$，则 $T(x) = Ax + b \pmod 1$，其中 $b = (\alpha, 0)^T$。
    
    \textbf{1. 保测性：}
    $\det(A) = 5(1) - 2(2) = 1$。
    由于 $A \in SL(2, \mathbb{Z})$，线性部分保测，平移部分也保测。因此 $T$ 是\textbf{保测}的。

    \textbf{2. 遍历性：}
    这是一个仿射变换。对于仿射变换 $Tx = Ax + b$，如果线性部分 $A$ 是遍历的（即无单位根特征值），且 $A$ 没有特征值 1，则 $T$ 与 $A$ 拓扑共轭（通过平移），从而也是遍历的。
    
    检查 $A$ 的特征值：
    $\lambda^2 - \text{tr}(A)\lambda + \det(A) = \lambda^2 - 6\lambda + 1 = 0$.
    $\lambda = \frac{6 \pm \sqrt{36-4}}{2} = 3 \pm \sqrt{8}$.
    特征值 $\lambda_1 = 3 + 2\sqrt{2} > 1$, $\lambda_2 = 3 - 2\sqrt{2} \in (0, 1)$。
    特征值均为无理数，且模不为 1，不是单位根。因此线性部分 $\sigma_A$ 是遍历的（甚至是 Anosov 的）。
    
    由于 $1$ 不是 $A$ 的特征值，方程 $(A-I)v = -b$ 有唯一解 $v = -(A-I)^{-1}b$。
    令 $S(x) = x - v$，则 $S \circ T \circ S^{-1}(y) = S(T(y+v)) = A(y+v) + b - v = Ay + Av + b - v = Ay + (A-I)v + b = Ay - b + b = Ay$。
    即 $T$ 与线性自同构 $\sigma_A$ 共轭。
    因为 $\sigma_A$ 是遍历的，所以 $T$ 也是\textbf{遍历}的。
    
    \textbf{结论}：$T$ 是保测且遍历的映射。
\end{solution}

\begin{problem}[10分]
    设多项式 $P(n) = \alpha n^d + a_{d-1}n^{d-1} + \dots + a_0$，其中 $\alpha \in \mathbb{R} \setminus \mathbb{Q}$，$a_0, \dots, a_{d-1} \in \mathbb{R}$。
    证明：$\{P(n) \pmod 1 \mid n \in \mathbb{N}\}$ 在 $[0, 1]$ 区间中均匀分布。
\end{problem}

\begin{solution}
    利用 \textbf{Weyl 判别法}：序列 $\{x_n\}$ 在 $[0,1]$ 上均匀分布当且仅当对于任意非零整数 $h \in \mathbb{Z} \setminus \{0\}$，有
    $$
    \lim_{N \to \infty} \frac{1}{N} \sum_{n=1}^N e^{2\pi i h x_n} = 0
    $$
    
    我们要证明对于 $x_n = P(n)$，上述极限为 0。
    这可以使用 \textbf{van der Corput 差分定理}（或 Weyl 的多项式定理）。
    
    \textbf{定理}：设 $\{x_n\}$ 是实数列。如果对于任意正整数 $k$，差分序列 $\Delta_k x_n = x_{n+k} - x_n$ 模 1 均匀分布，则 $\{x_n\}$ 模 1 均匀分布。
    
    \textbf{归纳法证明}：
    对多项式的次数 $d$ 进行归纳。
    1. 当 $d=1$ 时，$P(n) = \alpha n + a_0$。由于 $\alpha$ 是无理数，这是经典的无理旋转，已知是均匀分布的。
    2. 假设对于次数为 $d-1$ 的首项系数为无理数的多项式结论成立。
    3. 考虑次数为 $d$ 的多项式 $P(n) = \alpha n^d + \dots$。
       考察差分 $\Delta_k P(n) = P(n+k) - P(n)$。
       $$
       P(n+k) - P(n) = \alpha(n+k)^d - \alpha n^d + \dots = \alpha(n^d + d k n^{d-1} + \dots) - \alpha n^d + \dots
       $$
       $$
       = \alpha d k n^{d-1} + O(n^{d-2})
       $$
       这是一个关于 $n$ 的 $d-1$ 次多项式。
       其首项系数为 $\alpha d k$。
       由于 $\alpha$ 是无理数，$d, k$ 是非零整数，所以 $\alpha d k$ 也是无理数。
       根据归纳假设，对于任意 $k \ge 1$，序列 $\{\Delta_k P(n)\}$ 是均匀分布的。
       由 van der Corput 定理（或 Weyl 归纳步骤），原序列 $\{P(n)\}$ 也是均匀分布的。
       
    \textbf{证毕}。
\end{solution}

\begin{problem}[10分]
    设 $l \in \mathbb{N}_{>0}$。定义数列 $p_{n+1} = l \cdot p_n + p_{n-1}$，$p_0=0, p_1=1$。证明对于几乎所有的实数 $\alpha$，$\{p_n \cdot \alpha \pmod 1 \mid n \in \mathbb{N}\}$ 在 $[0, 1]$ 中均匀分布。
\end{problem}

\begin{solution}
    数列 $p_n$ 的递推关系对应于矩阵 $M = \begin{pmatrix} l & 1 \\ 1 & 0 \end{pmatrix}$ 的幂次。
    特征方程为 $\lambda^2 - l\lambda - 1 = 0$。
    特征值为 $\lambda_1 = \frac{l + \sqrt{l^2+4}}{2} > 1$ 和 $\lambda_2 = \frac{l - \sqrt{l^2+4}}{2}$。
    注意 $\lambda_1 \lambda_2 = -1$，且 $|\lambda_2| < 1$（因为 $l \ge 1$）。
    
    通项公式为 $p_n = c_1 \lambda_1^n + c_2 \lambda_2^n$。
    由于 $p_n$ 是整数，且 $\lambda_2^n \to 0$，所以 $p_n$ 表现为指数增长序列，类似于 $c \lambda_1^n$。
    
    这是一个关于 lacunary sequence（缺项级数/指数增长序列）的均匀分布问题。
    \textbf{定理 (Koksma)}：设 $\{u_n\}$ 是整数序列且 $|u_n| \ge 1$。如果存在 $\theta > 1$ 使得 $|u_{n+1}| \ge \theta |u_n|$ 对足够大的 $n$ 成立（即序列按指数增长），则对于几乎所有的 $\alpha \in \mathbb{R}$，序列 $\{u_n \alpha\}$ 模 1 均匀分布。
    
    在本题中，$p_{n+1} \approx \lambda_1 p_n$，且 $\lambda_1 > 1$。
    具体来说，$p_{n+1} = l p_n + p_{n-1} \ge 1 \cdot p_n + 0 = p_n$（单调递增）。
    实际上 $\lim p_{n+1}/p_n = \lambda_1 > 1$。
    因此满足 Koksma 定理的条件。
    所以对于几乎所有的 $\alpha$，$\{p_n \alpha\}$ 是均匀分布的。
\end{solution}

\begin{problem}[10分]
    构造紧致度量空间上的连续映射 $T: X \to X$，使得其拓扑熵等于 $\ln(l \cdot (\sqrt{5}+1))$（注：根据试卷上下文及上一题，此处应为 $\ln(\lambda_1) = \ln(\frac{l+\sqrt{l^2+4}}{2})$，或者若 $l=1$ 则是 $\ln(\frac{1+\sqrt{5}}{2})$。题目文字可能是 $\ln(\frac{l+\sqrt{l^2+4}}{2})$）。
    
    \textit{假设题目意图是构造熵为上一题矩阵特征值的映射。}
\end{problem}

\begin{solution}
    取 $X = \mathbb{T}^2 = \mathbb{R}^2/\mathbb{Z}^2$ 为二维环面。
    定义 $T: \mathbb{T}^2 \to \mathbb{T}^2$ 为由矩阵 $A = \begin{pmatrix} l & 1 \\ 1 & 0 \end{pmatrix}$ 诱导的线性自同构。
    即 $T\begin{pmatrix} x \\ y \end{pmatrix} = \begin{pmatrix} l & 1 \\ 1 & 0 \end{pmatrix} \begin{pmatrix} x \\ y \end{pmatrix} \pmod 1$。
    
    该映射的拓扑熵由 Bowen 公式给出：
    $$
    h(T) = \sum_{|\lambda| > 1} \ln |\lambda|
    $$
    其中 $\lambda$ 是 $A$ 的特征值。
    $A$ 的特征值为 $\lambda_{1,2} = \frac{l \pm \sqrt{l^2+4}}{2}$。
    其中 $\lambda_1 = \frac{l + \sqrt{l^2+4}}{2} > 1$。
    $\lambda_2 = \frac{l - \sqrt{l^2+4}}{2}$，其绝对值 $|\lambda_2| = |\frac{-1}{\lambda_1}| < 1$。
    因此，
    $$
    h(T) = \ln(\lambda_1) = \ln\left(\frac{l + \sqrt{l^2+4}}{2}\right)
    $$
    这给出了所需的构造。
\end{solution}

\begin{problem}[15分]
    设连续映射 $T: [0, 1] \to [0, 1]$ 为 $x \mapsto x^3$。试求 $T$ 的所有极小系统和不变概率测度。
\end{problem}

\begin{solution}
    \textbf{1. 求所有极小系统：}
    
    极小系统是指非空闭不变集 $M \subseteq [0, 1]$，且 $M$ 中不包含真闭不变子集。
    
    考察映射 $T(x) = x^3$ 的动力学性质：
    \begin{itemize}
        \item \textbf{不动点}：解方程 $x^3 = x$，在 $[0, 1]$ 上得到不动点 $x=0$ 和 $x=1$。
        \item \textbf{轨道分析}：
        \begin{itemize}
            \item 若 $x=0$，则 $T(0)=0$。$\{0\}$ 是闭不变集。
            \item 若 $x=1$，则 $T(1)=1$。$\{1\}$ 是闭不变集。
            \item 若 $x \in (0, 1)$，由于 $x^3 < x$，轨道 $\{T^n(x)\}$ 单调递减。由单调有界原理，$\lim_{n\to\infty} T^n(x) = 0$。
        \end{itemize}
    \end{itemize}

    假设 $M$ 是一个极小集：
    \begin{itemize}
        \item 显然 $\{0\}$ 和 $\{1\}$ 都是极小集。
        \item 若 $M$ 包含 $(0, 1)$ 中的点 $x_0$，由于 $M$ 是不变集，则整个轨道 $\{T^n(x_0)\} \subseteq M$。
        \item 由于 $M$ 是闭集，轨道的极限点 $0 \in M$。
        \item 这意味着 $\{0\} \subseteq M$。根据极小性的定义，必须有 $M = \{0\}$。这与 $x_0 \in M$ 且 $x_0 \neq 0$ 矛盾。
    \end{itemize}
    因此，不存在包含 $(0, 1)$ 中点的极小集。
    
    \textbf{结论}：$T$ 的所有极小系统为 $\{0\}$ 和 $\{1\}$。

    \textbf{2. 求所有不变概率测度：}
    
    设 $\mu$ 是 $[0, 1]$ 上的 Borel 概率测度。$\mu$ 是 $T$-不变的，即对任意 Borel 集 $A$，有 $\mu(T^{-1}(A)) = \mu(A)$。
    
    根据 Krylov-Bogolyubov 定理，不变测度的支撑集 $\text{supp}(\mu)$ 包含在非游荡集中。对于本题，我们可以直接分析测度的分布。
    
    考虑开区间 $(0, 1)$。对于任意 $0 < a < 1$，令 $A = (a, 1]$。
    $$ T^{-1}(A) = \{x \in [0, 1] \mid x^3 > a\} = (a^{1/3}, 1] $$
    由不变性，$\mu((a, 1]) = \mu((a^{1/3}, 1])$。
    
    令 $F(x) = \mu((x, 1])$。则有 $F(a) = F(a^{1/3})$。
    对任意 $a \in (0, 1)$，迭代该关系：
    $$ F(a) = F(a^{1/3}) = F(a^{1/9}) = \dots = F(a^{1/3^n}) $$
    当 $n \to \infty$ 时，$a^{1/3^n} \to 1$。
    利用测度的连续性（或分布函数的性质）：
    $$ F(a) = \lim_{n \to \infty} F(a^{1/3^n}) = \mu(\{1\}) $$
    这意味着对于任意 $a \in (0, 1)$，$\mu((a, 1]) = \mu(\{1\})$。
    从而 $\mu((a, 1)) = \mu((a, 1]) - \mu(\{1\}) = 0$。
    令 $a \to 0$，由测度的连续性得 $\mu((0, 1)) = 0$。
    
    因此，任何不变概率测度 $\mu$ 的质量只能集中在不动点 $0$ 和 $1$ 上。
    即 $\mu$ 必须是点测度 $\delta_0$ 和 $\delta_1$ 的线性组合。
    
    \textbf{结论}：$T$ 的所有不变概率测度为
    $$ \mu = p \delta_0 + (1-p) \delta_1, \quad \text{其中 } 0 \le p \le 1 $$
\end{solution}

\begin{problem}[15分]
    证明数列 $\{\sqrt{\ln n} \pmod 1 \mid n \in \mathbb{N}\}$ 在 $[0, 1]$ 区间中稠密但非均匀分布。（考虑 $\sqrt{\ln(n+1)} - \sqrt{\ln n}$）
\end{problem}

\begin{solution}
    令 $x_n = \sqrt{\ln n}$。
    
    \textbf{1. 证明稠密性：}
    
    考察相邻项的差分 $\Delta x_n = x_{n+1} - x_n$：
    $$
    \Delta x_n = \sqrt{\ln(n+1)} - \sqrt{\ln n} = \sqrt{\ln n + \ln(1+\frac{1}{n})} - \sqrt{\ln n}
    $$
    利用泰勒展开 $\ln(1+x) \approx x$ 和 $\sqrt{a+b} - \sqrt{a} = \frac{b}{\sqrt{a+b}+\sqrt{a}}$：
    $$
    \Delta x_n \approx \frac{\ln(1+1/n)}{\sqrt{\ln n} + \sqrt{\ln n}} \approx \frac{1/n}{2\sqrt{\ln n}} = \frac{1}{2n\sqrt{\ln n}}
    $$
    显然，当 $n \to \infty$ 时，$\Delta x_n \to 0$。
    同时，由于 $\lim_{n\to\infty} \ln n = +\infty$，故 $x_n \to +\infty$。
    
    \textbf{定理}：若实数列 $\{x_n\}$ 满足 $x_n \to \infty$ 且 $\Delta x_n \to 0$，则 $\{x_n \pmod 1\}$ 在 $[0, 1]$ 中稠密。
    
    \textit{直观解释}：由于步长趋于 0 且总行程无限，序列就像一只步子越来越小的蜗牛，必然会经过圆周上任意一点的任意邻域无数次。
    
    \textbf{2. 证明非均匀分布：}
    
    我们使用 Fejér 定理的一个推论，或者直接通过计数论证。
    
    \textbf{判别法}：设 $x_n = f(n)$，其中 $f(t)$ 是可微函数。若 $f'(t) \to 0$ 单调递减，则 $\{x_n\}$ 模 1 均匀分布的充要条件是 $\lim_{t\to\infty} t |f'(t)| = +\infty$。
    
    在本题中，$f(t) = \sqrt{\ln t}$，其导数为 $f'(t) = \frac{1}{2t\sqrt{\ln t}}$。
    考察极限：
    $$ \lim_{t\to\infty} t f'(t) = \lim_{t\to\infty} \frac{1}{2\sqrt{\ln t}} = 0 $$
    由于极限为 0（不为 $+\infty$），故该序列\textbf{不是}均匀分布的。
    
    \textbf{直接计数论证（更直观）}：
    考察 $x_n$ 的整数部分为 $k$ 时，其小数部分 $\{x_n\}$ 在 $[0, 1)$ 上的分布。
    $k \le \sqrt{\ln n} < k+1 \iff e^{k^2} \le n < e^{(k+1)^2}$。
    该区间内的整数个数为 $N_k \approx e^{(k+1)^2} - e^{k^2}$。
    
    将该区间分为前半段（小数部分在 $[0, 0.5)$）和后半段（小数部分在 $[0.5, 1)$）：
    \begin{itemize}
        \item 前半段 $n$ 的范围：$e^{k^2} \le n < e^{(k+0.5)^2}$。
        \item 后半段 $n$ 的范围：$e^{(k+0.5)^2} \le n < e^{(k+1)^2}$。
    \end{itemize}
    比较两段的长度（即落入点的个数）：
    $$ \frac{\text{后半段长度}}{\text{前半段长度}} \approx \frac{e^{(k+1)^2} - e^{(k+0.5)^2}}{e^{(k+0.5)^2} - e^{k^2}} \approx \frac{e^{k^2+2k}}{e^{k^2+k}} = e^k $$
    当 $k \to \infty$ 时，后半段的点数是前半段的 $e^k$ 倍。这意味着随着 $n$ 增大，绝大多数点都集中在区间的右端（即小数部分接近 1 的地方）。
    因此，该序列极度不均匀，概率测度弱收敛于 $\delta_1$（即 $\delta_0$ 在圆周上）。
\end{solution}

\section*{三、做题方法总结}

\begin{enumerate}
    \item \textbf{环面自同构的遍历性判定}：
    对于 $T(x) = Ax \pmod 1$（$A$ 为整数矩阵）：
    \begin{itemize}
        \item \textbf{保测性}：检查 $|\det(A)| = 1$。如果是，则保 Lebesgue 测度。
        \item \textbf{遍历性}：检查 $A$ 的特征值。如果 $A$ 没有单位根特征值（即 $\lambda^k \neq 1$），则 $T$ 是遍历的。对于 $SL(n, \mathbb{Z})$ 矩阵，这等价于没有模为 1 的特征值（双曲性），或者特征多项式在 $\mathbb{Z}[x]$ 上不可约且不是分圆多项式等更精细的条件。最简单的判据是：无特征值是单位根。
    \end{itemize}

    \item \textbf{拓扑熵的计算}：
    \begin{itemize}
        \item 对于环面自同构 $A$，拓扑熵 $h(T) = \sum_{|\lambda_i|>1} \ln |\lambda_i|$。
        \item 对于圆周映射，利用度数或共轭模型。
    \end{itemize}

    \item \textbf{均匀分布的判定 (Weyl 判别法)}：
    \begin{itemize}
        \item 证明 $\frac{1}{N} \sum e^{2\pi i h x_n} \to 0$。
        \item 利用 van der Corput 差分定理处理多项式形式。
        \item 利用度量理论（Koksma 定理）处理指数增长序列 $c \lambda^n$。
    \end{itemize}

    \item \textbf{区间映射的动力学分析}：
    \begin{itemize}
        \item 寻找不动点和周期点。
        \item 分析不动点间的单调性（如 $x^3 < x$）。
        \item 极小集通常是不动点或周期轨道（对于简单映射）。
        \item 不变测度通常支撑在非游荡集（极小集）上。
    \end{itemize}
\end{enumerate}

\end{document}
